\documentclass{beamer}

\usepackage{tikz}
\usetikzlibrary{positioning,matrix,fit, arrows.meta}
\usetikzlibrary{graphdrawing.trees}
\usepackage{color}
\newcommand{\vet}[1]{\foreach \num in {#1}{\el{\num}}}
\newcommand{\vetempty}[1]{\foreach \num in {#1}{\elempty{\num}}}
\newcommand{\el}[1]{\tikz{\node[font=\smaller, draw]{#1}}}
\newcommand{\elempty}[1]{\tikz{\node[font=\small, minimum size=1cm, draw=white, fill=white]{#1}}}
\newcommand{\bl}{\color{blue}}
\newcommand{\gr}{\color{gray}}
\newcommand{\re}{\color{red}}
\newcommand{\tcb}[1]{\textcolor{blue}{\textbf{#1}}}
\newcommand{\tcg}[1]{\textcolor{green}{\textbf{#1}}}
\newcommand{\tcr}[1]{\textcolor{red}{\textbf{#1}}}
\newcommand{\ar}{\tikz{\draw [->] (0,0) -- (0,0.4)}}
\newcommand{\nd}{\color{white!0}.}

\usepackage{beamerthemeSingapore}
\usepackage{graphicx}
%\usepackage{here}
\usepackage{url}
\usepackage[brazil]{babel}
\usepackage[T1]{fontenc}
\usepackage[utf8]{inputenc}
%\usepackage[tight,raggedright]{subfigure}
\usepackage{subfigure}
\usepackage[alf]{abntex2cite}
\usepackage{listings}
\usepackage{minted}
\usepackage{algorithm2e}
\usepackage{amsmath}
\usepackage{pdfpages}
% Please add the following required packages to your document preamble:
\usepackage{booktabs}
\usepackage{ragged2e}
\usepackage{etoolbox}
\usepackage{mathtools}
\usepackage{caption}
\usepackage{xmpmulti}

\tikzset{
  treenode/.style = {align=center, inner sep=0pt, text centered,
    font=\sffamily},
  arn_b/.style = {treenode, circle, black, font=\sffamily\bfseries, draw=black,
    fill=white, text width=1.5em},% arbre rouge noir, noeud noir
  arn_n/.style = {treenode, circle, white, font=\sffamily\bfseries, draw=black,
    fill=black, text width=1.5em},% arbre rouge noir, noeud noir
  arn_r/.style = {treenode, circle, red, draw=red, 
    text width=1.5em, very thick},% arbre rouge noir, noeud rouge
  arn_x/.style = {treenode, rectangle, draw=black,
    minimum width=0.5em, minimum height=0.5em}% arbre rouge noir, nil
}

\captionsetup[figure]{labelformat=empty,textfont=smaller}

\DeclarePairedDelimiter{\ceil}{\lceil}{\rceil}

\apptocmd{\frame}{}{\justifying}{} % Allow optional arguments after frame.

% TODO TODO TODO TODO TODO TODO TODO
\title{07 - Percursos e Propriedades de Árvores}
\author{Prof. Alex Kutzke}
\date{Agosto de 2025}

\begin{document}

% ==========================================
\begin{frame}
  \begin{center}
    % TODO TODO TODO TODO TODO TODO
    \large{Setor de Educação Profissional e Tecnológica}\\
    Tecnologia em Análise e Desenvolvimento de Sistemas\\
    \vspace{1cm}			
    \small{DS141 - Estruturas de Dados II}
  \end{center}
  \titlepage

  \begin{center}
    Slides adaptados do material produzido\\pelo Prof. Rodrigo Minetto
  \end{center}
\end{frame}
% ==========================================

% ==========================================
\begin{frame}[c]{Roteiro}
  \tableofcontents

\end{frame}
% ==========================================

\section{Percursos}

\begin{frame}[c]{Percurso}
  \begin{itemize}
    \item Um \textbf{percurso} corresponde a uma \tcr{visita} sistemática a cada um dos
      nós da árvore;
    \item Muitas operações em árvores binárias envolvem o percurso de todas as subárvores,
      con a execução de alguma ação em cada nó.
  \end{itemize}
\end{frame}

\begin{frame}[c]{Percurso}
  \begin{itemize}
    \item Por exemplo: se cada nó da árvore possui um campo que armazena o salário;
    \item Para realizar um ``reajuste salarial'', precisamos alterar o valor de todos os nós;
    \item Nesse caso ``visitar um nó'' seria atualizar o valor armazenado;
    \item Não podemos esquecer de nenhum nó, nem queremos visitar um nó mais do que uma vez;
    \item \textbf{Nesse caso}, a ordem da visita não é importante:
    \begin{itemize}
      \item Mas em algumas outras aplicações, queremos visitar os nós em certa ordem desejada.
    \end{itemize}
  \end{itemize}
\end{frame}

\begin{frame}
  \frametitle{Visita}

  \textbf{Visitar} um nó significa trabalhar com
a informação do nó (por exemplo: imprimir, modificar, etc). Contudo, 
durante um percurso podemos passar várias vezes por alguns dos nós sem
visitá-los.
\end{frame}


\begin{frame}
  \frametitle{Formas de Percurso}

  É comum percorrer uma árvore em uma das seguintes ordens:
  \begin{itemize}
    \item Pré-ordem (\tcb{R,E,D});
    \item In-ordem (\tcr{E,R,D});
    \item Pós-ordem (\tcg{E,D,R});
  \end{itemize}

\end{frame}

\begin{frame}
  \frametitle{Pré-ordem}

  \begin{description}
    \item [Pré-ordem] (R, E, D):
    \begin{enumerate}
      \item Visitar a \textbf{raiz};
      \item Percorrer a \textbf{subárvore da esquerda} em pré-ordem;
      \item Percorrer a \textbf{subárvore da direita} em pré-ordem.
    \end{enumerate}
  \end{description}
  
\end{frame}

\begin{frame}
  \frametitle{Pré-ordem: exemplo}

  Percurso pré-ordem para a árvore abaixo:\\
  \texttt{50, 30, 20, 10, 40, 35, 45, 90, 95}
  
  \vspace{1cm}

  \begin{columns}
    \begin{column}{0.6\textwidth}

      \centering
      \begin{tikzpicture}[-]
        \node (root) [arn_b] {50}
        child { 
        node [arn_b] {30} 
        child { 
        node [arn_b] {20} 
        child { node [arn_b] {10} }
        child [missing]
        } 
        child { 
        node [arn_b] {40}
        child { node [arn_b] {35} }
        child { node [arn_b] {45} }
        } 
        }
        child [missing]
        child { 
        node [arn_b] {90} 
        child [missing]
        child { node [arn_b] {95} } 
        };
      \end{tikzpicture}
    \end{column}
    \begin{column}{0.4\textwidth}
      \begin{itemize}
        \item Segue os nós até chegar ao mais \tcr{"profundos"}, em \tcr{"ramos"} de subárvores da esquerda para a direita;
        \item Percurso em \tcb{profundidade}.
      \end{itemize}
    \end{column}
  \end{columns}

\end{frame}

\begin{frame}
  \frametitle{In-ordem}

  \begin{description}
    \item [In-ordem] (E, R, D):
    \begin{enumerate}
      \item Percorrer a \textbf{subárvore da esquerda} em in-ordem;
      \item Visitar a \textbf{raiz};
      \item Percorrer a \textbf{subárvore da direita} em in-ordem.
    \end{enumerate}
  \end{description}
  
\end{frame}

\begin{frame}
  \frametitle{In-ordem: exemplo}

  Percurso in-ordem para a árvore abaixo:\\
  \texttt{10, 20, 30, 35, 40, 45, 50, 90, 95}
  
  \vspace{1cm}

  \begin{columns}
    \begin{column}{0.6\textwidth}

      \centering
      \begin{tikzpicture}[-]
        \node (root) [arn_b] {50}
        child { 
        node [arn_b] {30} 
        child { 
        node [arn_b] {20} 
        child { node [arn_b] {10} }
        child [missing]
        } 
        child { 
        node [arn_b] {40}
        child { node [arn_b] {35} }
        child { node [arn_b] {45} }
        } 
        }
        child [missing]
        child { 
        node [arn_b] {90} 
        child [missing]
        child { node [arn_b] {95} } 
        };
      \end{tikzpicture}
    \end{column}
    \begin{column}{0.4\textwidth}
      \begin{itemize}
        \item Visita a raiz entre as ações de percorrer as duas subárvores;
        \item Percurso em \tcb{ordem simétrica}.
      \end{itemize}
    \end{column}
  \end{columns}

\end{frame}


\begin{frame}
  \frametitle{Pós-ordem}

  \begin{description}
    \item [Pós-ordem] (E, D, R):
    \begin{enumerate}
      \item Percorrer a \textbf{subárvore da esquerda} em pós-ordem;
      \item Percorrer a \textbf{subárvore da direita} em pós-ordem.
      \item Visitar a \textbf{raiz};
    \end{enumerate}
  \end{description}
  
\end{frame}

\begin{frame}
  \frametitle{Pós-ordem: exemplo}

  Percurso pré-ordem para a árvore abaixo:\\
  \texttt{10, 20, 35, 45, 40, 30, 95, 90, 50}
  
  \vspace{1cm}

  \begin{columns}
    \begin{column}{0.6\textwidth}

      \centering
      \begin{tikzpicture}[-]
        \node (root) [arn_b] {50}
        child { 
        node [arn_b] {30} 
        child { 
        node [arn_b] {20} 
        child { node [arn_b] {10} }
        child [missing]
        } 
        child { 
        node [arn_b] {40}
        child { node [arn_b] {35} }
        child { node [arn_b] {45} }
        } 
        }
        child [missing]
        child { 
        node [arn_b] {90} 
        child [missing]
        child { node [arn_b] {95} } 
        };
      \end{tikzpicture}
    \end{column}
    \begin{column}{0.4\textwidth}
      \begin{itemize}
        \item Visita a raiz \textbf{após} as ações de percorrer as duas subárvores;
      \end{itemize}
    \end{column}
  \end{columns}

\end{frame}

\section{Altura da árvore}

\begin{frame}
  \frametitle{Altura da árvore}

  \begin{itemize}
    \item Existe apenas um caminho da raiz para qualquer nó:
      \begin{itemize}
        \item Em outras palavras, o que isso quer dizer?;
        \pause
        \item Não existem ligações entre nós ``irmãos''. Quando isso ocorre,
              a estrutura deixa de ser uma árvore de se torna um \textit{grafo};
      \end{itemize}
    \pause
    \item Assim, \tcb{altura} é:

    \center \textit{o comprimento do caminho mais longo, da raiz até uma das folhas.}

  \end{itemize}
\end{frame}

\begin{frame}
  \frametitle{Altura da árvore}

  \begin{itemize}
    \item Altura de uma árvore com um único nó é 0;
    \item Altura de uma árvore vazia é -1;
  \end{itemize}

  \begin{columns}[t]
    \begin{column}{0.3\textwidth}
      \centering
      Altura 0\\
      \begin{tikzpicture}[-]
        \node (root) [arn_b] {};
      \end{tikzpicture}
    \end{column}
    \begin{column}{0.3\textwidth}
      \centering
      Altura 1\\
      \begin{tikzpicture}[-]
        \node (root) [arn_b] {1}
        child [missing] 
        child { node [arn_b] {2} };
      \end{tikzpicture}
    \end{column}
    \begin{column}{0.3\textwidth}
      \centering
      Altura 2\\
      \begin{tikzpicture}[-]
        \node (root) [arn_b] {}
        child { 
          node [arn_b] {} 
        }
        child [missing]
        child { 
          node [arn_b] {} 
          child { node [arn_b] {} }
          child [missing]
        };
      \end{tikzpicture}
    \end{column}
  \end{columns}
\end{frame}


\begin{frame}
  \frametitle{Níveis}

  \centering
  \begin{tikzpicture}[-]
    \node (root) [arn_b] {}
    child { 
      node [arn_b] {} 
    }
    child [missing]
    child { 
      node (lv1) [arn_b] {} 
      child { node (lv2) [arn_b] {} }
      child [missing]
    };
    \node [fit=(root)(lv1)(lv2)] (box) {};
    \node (lroot) [right = 3cm of root] {Nível 0};
    \node (llv1) [below= of lroot]  {Nível 1};
    \node (llv2) [below= of llv1]  {Nível 2};
    \draw[->,shorten <= 5pt,dashed] (root) -- (lroot);
    \draw[->,shorten <= 5pt,dashed] (lv1) -- (llv1);
    \draw[->,shorten <= 5pt,dashed] (lv2) -- (llv2);
  \end{tikzpicture}
\end{frame}

\section{Altura, número de nós e eficiência}

\begin{frame}
  \frametitle{Árvores cheias: número de nós}

    \only<1>{\center Uma árvore \textit{binária} \tcb{cheia (ou completa)} é aquela em que todos os nós internos têm duas subárvores e todos os nós folhas estão no último nível.}

    \only<2>{\center Podemos notar que nesse tipo de árvore temos \textbf{um} nó no nível \tcb{0}, \text{dois} nós no nível \tcb{1}, \textbf{quatro} nós no nível \tcb{2}, \textbf{oito} nós no nível \tcb{3}, e assim por diante. Isto é, no nível $n$, temos $2^n$ nós.}

    \only<3>{\center É possível mostrar que uma árvore cheia de altura $h$ tem um número de nós dado por \textbf{$2^{h+1} - 1$}. \tcb{Como?}}

    \vspace{1cm}

    \centering
    \begin{tikzpicture}[-]
      \node (root) [arn_b] {}
      child { 
        node [arn_b] {} 
          child { 
            node [arn_b] {} 
          }
          child { 
            node [arn_b] {} 
          }
      }
      child [missing]
      child { 
        node [arn_b] {} 
          child { 
            node [arn_b] {} 
          }
          child { 
            node [arn_b] {} 
          }
      };
    \end{tikzpicture}
\end{frame}

\begin{frame}
  \frametitle{Árvores degeneradas}

  \begin{itemize}
    \item Uma árvore é dita \tcb{degenerada} se todos os seus
      nós internos tem uma única subárvore associada;
    \item Apenas um nó por nível;
    \item Uma árvore degenerada de altura $h$ possui $h+1$ nós;
  \end{itemize}

  \vspace{1cm}

  \begin{columns}
    \begin{column}{0.5\textwidth}
      \centering
      \begin{tikzpicture}[-,level distance=1cm]
        \node (root) [arn_b] {}
          child { 
            node [arn_b] {}
            child { 
              node [arn_b] {} 
              child { 
                node [arn_b] {} 
              }
              child [missing]
            }
            child [missing]
          }
          child [missing];
      \end{tikzpicture}
    \end{column}
    \begin{column}{0.5\textwidth}
      \centering
      \begin{tikzpicture}[-, level distance=1cm]
        \node (root) [arn_b] {}
          child { 
            node [arn_b] {}
            child [missing]
            child { 
              node [arn_b] {} 
              child { 
                node [arn_b] {} 
              }
              child [missing]
            }
          }
          child [missing];
      \end{tikzpicture}
    \end{column}
  \end{columns}

\end{frame}

\begin{frame}[c]
  \frametitle{Eficiência X Altura}

  \centering
  A altura de uma árvore é uma medida importante 
  na avaliação da eficiência com que visitamos os 
  nós da árvore. 

  \vspace{1cm}

  Uma árvore binária com $n$ nós tem uma altura mínima
  proporcional a $\log_2{n}$ (caso da árvore cheia) e
  uma altura máxima proporcional a $n$ (caso da árvore 
  degenerada). A altura indica o esforço computacional 
  necessário para alcançar qualquer nó da árvore.
\end{frame}

\section{Exercícios propostos}

\begin{frame}
  \frametitle{Exercícios propostos}

  \textbf{Exercício 1} Suponha a árvore binária abaixo. (a) ache o grau
  de cada um dos nós; (b) determine os nós folhas; e (c) complete a 
  árvore com quantos nós forem necessários para transformá-la em uma
  árvore binária cheia.
      
  \centering
  \begin{tikzpicture}[-, level distance=1cm]
    \node (root) [arn_b] {50}
    child { 
      node [arn_b] {30} 
      child { 
        node [arn_b] {20} 
        child { node [arn_b] {10} }
        child [missing]
      } 
      child { 
        node [arn_b] {40}
        child [missing]
        child [missing]
      } 
    }
    child [missing]
    child { 
      node [arn_b] {90} 
      child [missing]
      child { 
        node [arn_b] {95} 
        child [missing]
        child { 
          node [arn_b] {98} 
        } 
      } 
    };
  \end{tikzpicture}
\end{frame}

\begin{frame}
  \frametitle{Exercícios propostos}

  \textbf{Exercício 2} Escreva funções que recebam uma árvore binária
  inteira como entrada e a imprima com os seguintes tipos de percurso:

  \begin{itemize}
    \item pré-ordem;
    \item in-ordem;
    \item pós-ordem.
  \end{itemize}
\end{frame}

\begin{frame}
  \frametitle{Exercícios propostos}

  \textbf{Exercício 3} Escreva uma função que conte o número de nós
  de uma árvore binária. Utilize o seguinte protótipo para a sua função:

  \texttt{int conta\_nos(Arvore \*a);}
  
  \vspace{1cm}

  \textbf{Exercício 4} Escreva uma função que imprime e retorna o maior
  valor inteiro em uma árvore binária composta por valores inteiros. 
  Utilize o seguinte protótipo para a sua função:

  \texttt{int max\_arvore(Arvore \*a);}

\end{frame}

\begin{frame}
  \frametitle{Exercícios propostos}

  \textbf{Exercício 5} Escreva uma função que calcula a altura de uma árvore
  binária. Utilize o seguinte protótipo para a sua função:

  \texttt{int altura\_arvore(Arvore *a);}

  \vspace{1cm}

  \textbf{Exercício 6} Escreva uma função que conta o número de nós
  folhas em uma árvore binária.
  Utilize o seguinte protótipo para a sua função:

  \texttt{int nos\_folha\_arvore(Arvore *a);}

\end{frame}

% ==========================================
% Frame para referências
\begin{frame}[allowframebreaks]
  \frametitle{Referências utilizadas na apresentação}
  \bibliography{referencias}
\end{frame}
% ==========================================

\end{document}
